\documentclass[11pt, a4paper]{article}
\usepackage[utf8]{inputenc}
\usepackage{amsmath, amssymb}
\usepackage{geometry}
\usepackage{parskip}
\usepackage{enumitem}
\usepackage{booktabs}
\usepackage{listings}
\usepackage{xcolor}
\usepackage{mdframed}
\usepackage{hyperref}

\geometry{top=2.5cm, bottom=2.5cm, left=2.8cm, right=2.8cm}

\definecolor{enunbg}{RGB}{240,240,250}
\definecolor{enunline}{RGB}{80,80,150}
\definecolor{codegray}{RGB}{250,250,250}
\definecolor{codeblue}{RGB}{0,80,160}
\definecolor{codegreen}{RGB}{0,120,0}
\definecolor{codecomment}{RGB}{100,100,100}

\lstdefinestyle{pystyle}{
  language=Python, backgroundcolor=\color{codegray},
  basicstyle=\footnotesize\ttfamily, commentstyle=\color{codecomment},
  keywordstyle=\color{codeblue}\bfseries, stringstyle=\color{codegreen},
  numberstyle=\tiny\color{gray}, breaklines=true,
  numbers=left, numbersep=5pt, showstringspaces=false, tabsize=4,
  frame=single, rulecolor=\color{gray!40}, xleftmargin=12pt,
}

\newmdenv[backgroundcolor=enunbg, linecolor=enunline, linewidth=1pt,
  innerleftmargin=12pt, innerrightmargin=12pt, innertopmargin=8pt,
  innerbottommargin=8pt, skipabove=8pt, skipbelow=8pt]{enun}

\newcommand{\res}{\smallskip\noindent\textbf{\color{enunline}Resoluci\'on:}\smallskip}

\title{\textbf{Gu\'ia G5 --- Ventanas y Estimaci\'on Espectral}\\[0.2cm]
       \large An\'alisis y Procesamiento de Se\~nales --- UNSAM}
\date{}

\begin{document}
\maketitle

\section*{Antes de arrancar}

Esta gu\'ia cierra todo lo que vimos sobre an\'alisis espectral en se\~nales finitas:
ventanas, periodograma, Bartlett, Welch.
Los ejercicios alternan preguntas conceptuales con pr\'actica num\'erica.
Las referencias al Holton son al cap\'itulo 14.

% =========================================================
\section*{Ejercicio 1: ¿Cu\'ando aparece desparramo espectral?}

\begin{enun}
\textbf{a)} El desparramo espectral aparece en \emph{dos} situaciones distintas
al calcular la DFT de un bloque finito de se\~nal.
Indique cu\'ales son las dos situaciones.

\textbf{b)} Mire la Figura~14.6 del Holton (las cuatro ventanas: rectangular,
Hamming, Hann y Blackman).
?`Qu\'e dos par\'ametros de la ventana se ven en esa figura?
?`Cu\'al de ellos est\'a relacionado con cada una de las dos situaciones del punto a)?

\textbf{c)} Para $x[n] = \cos(2\pi \cdot 205 \cdot n/f_s)$ con $f_s = 1000$~Hz
y $N = 100$ muestras, calcule:
\begin{itemize}[noitemsep]
  \item El espaciado de bins $\Delta f = f_s / N$.
  \item El bin te\'orico al que corresponde 205~Hz ($k = f_0 \cdot N / f_s$).
  \item ?`Cae justo en un bin entero?
\end{itemize}
Corra el c\'odigo del final del ejercicio y mire el espectro: ?`d\'onde aparece
energ\'ia adem\'as del pico?
\end{enun}

\res

\textbf{a) Las dos situaciones.}
El desparramo (Figura~14.14 del Holton) aparece cuando:

\begin{enumerate}
  \item \textbf{Dos frecuencias caen dentro del l\'obulo principal} de la ventana.
    Como cada componente del espectro real est\'a convolucionada con la
    transformada de la ventana, los picos se ``ensanchan''.
    Si dos picos est\'an muy cerca, sus l\'obulos principales se superponen y
    aparece un solo pico en lugar de dos. Eso es p\'erdida de \emph{resoluci\'on}.

  \item \textbf{Una frecuencia no cae exactamente en un bin de la DFT.}
    La DFT solo eval\'ua el espectro en los bins discretos $k \cdot f_s/N$.
    Si la frecuencia real est\'a entre dos bins, el pico ``se reparte'' entre
    los bins vecinos y adem\'as aparece energ\'ia espuria en los dem\'as bins
    (los l\'obulos laterales de la ventana se ven). Eso es \emph{leakage}
    (filtraci\'on).
\end{enumerate}

\textbf{b) Par\'ametros de la ventana.}
En la Figura~14.6 se ven dos cosas para cada ventana:

\begin{itemize}
  \item El \textbf{ancho del l\'obulo principal} ($\Delta\omega$): cu\'an
    estrecho o ancho es el pico central de la transformada de la ventana.
    Esto controla la situaci\'on (1): l\'obulo principal m\'as estrecho
    $\Rightarrow$ mejor resoluci\'on.

  \item El \textbf{nivel del primer l\'obulo lateral}
    (cu\'antos dB por debajo del pico central est\'a el primer ``bulto''
    a los costados). Esto controla la situaci\'on (2):
    l\'obulos laterales m\'as bajos $\Rightarrow$ menos energ\'ia espuria.
\end{itemize}

\textbf{c) C\'alculo num\'erico.}
\begin{itemize}
  \item $\Delta f = 1000/100 = 10$~Hz por bin.
  \item $k = 205 \cdot 100 / 1000 = 20{,}5$.
  \item \textbf{No cae en un bin entero}: cae exactamente \emph{a mitad de
    camino} entre el bin 20 (200~Hz) y el bin 21 (210~Hz).
\end{itemize}

Al correr el c\'odigo, se ve que con \textbf{ventana rectangular} (sin enventanar)
aparecen ``patas'' que se extienden por todo el espectro, decayendo lentamente.
Con \textbf{ventana de Hann}, el pico se ve m\'as ancho pero las patas
desaparecen casi por completo. Eso muestra los dos efectos del enventanado:
suaviza el desparramo a costa de ensanchar el pico.

\begin{lstlisting}[style=pystyle, caption={Visualizaci\'on r\'apida del desparramo}]
import numpy as np
import matplotlib.pyplot as plt
import scipy.signal.windows as spsw

fs = 1000
N  = 100
n  = np.arange(N)
x  = np.cos(2*np.pi * 205 * n / fs)

freqs = np.fft.rfftfreq(N, d=1/fs)

# Sin ventana (rectangular) vs Hann
for nombre, w in [("Rectangular", np.ones(N)),
                  ("Hann", spsw.hann(N, sym=False))]:
    X = np.fft.rfft(x * w) / N
    plt.figure(figsize=(8, 3))
    plt.plot(freqs, 20*np.log10(np.abs(X) + 1e-12), label=nombre)
    plt.axvline(205, color="red", linestyle="--", label="205 Hz real")
    plt.xlabel("Frecuencia [Hz]")
    plt.ylabel("|X(f)| [dB]")
    plt.title(f"Espectro con ventana {nombre}")
    plt.ylim(-80, 0)
    plt.grid(True); plt.legend(); plt.tight_layout()
plt.show()
\end{lstlisting}

% =========================================================
\section*{Ejercicio 2: Tabla de ventanas}

\begin{enun}
La siguiente tabla (Holton, Sec.~14.1.2) resume las ventanas m\'as usadas:

\begin{center}
\small
\begin{tabular}{lccl}
\toprule
\textbf{Ventana} & \textbf{Ancho l\'ob.\ ppal.} & \textbf{1\textsuperscript{er} l\'ob.\ lat.} & \textbf{Para qu\'e sirve} \\
\midrule
Rectangular     & $4\pi/N$   & $-13$~dB  & m\'axima resoluci\'on \\
Hann            & $8\pi/N$   & $-32$~dB  & uso general \\
Hamming         & $8\pi/N$   & $-43$~dB  & se\~nales d\'ebiles cercanas \\
Blackman-Harris & $12\pi/N$  & $-92$~dB  & se\~nales muy d\'ebiles \\
Flattop         & $\approx 18\pi/N$ & $-93$~dB & medici\'on de amplitud \\
\bottomrule
\end{tabular}
\end{center}

\textbf{Definiciones que vamos a usar:}
\begin{itemize}[noitemsep]
  \item \textbf{Ancho del l\'obulo principal}: ancho en frecuencia del pico
    central de la ventana. M\'as estrecho $=$ mejor resoluci\'on.
  \item \textbf{Nivel del primer l\'obulo lateral} (en dB): qu\'e tan abajo
    est\'a el primer ``bulto'' a los costados del pico, medido respecto a la
    cima del pico (que est\'a en 0~dB). Cuanto m\'as negativo, mejor.
  \item \textbf{Enmascarar}: que el ``piso'' de los l\'obulos laterales de
    una se\~nal fuerte tape (no permita ver) a una se\~nal d\'ebil de otra
    frecuencia.
\end{itemize}

\textbf{a)} Para $f_s = 8000$~Hz y $N = 800$, calcule la resoluci\'on
m\'inima en Hz para cada ventana. Use:
\[
  \Delta f_{\min} = \frac{\Delta\omega}{2\pi} \cdot f_s
\]
?`Cu\'al de las ventanas permite separar dos tonos a $50$~Hz de distancia?

\textbf{b)} Tengo una se\~nal con un pico fuerte en $0$~dB y un pico d\'ebil
en $-40$~dB. ?`Qu\'e ventanas \emph{no enmascaran} al pico d\'ebil? (Use la
columna del primer l\'obulo lateral).

\textbf{c)} La ventana \textbf{Flattop} tiene el l\'obulo principal m\'as ancho
de todas, lo cual deber\'ia ser malo. Sin embargo se usa mucho en metrolog\'ia.
?`Por qu\'e? ?`Qu\'e propiedad de su l\'obulo principal la hace especial?
\end{enun}

\res

\textbf{a) Resoluci\'on m\'inima.}
El espaciado de bins base es $\Delta f_{\text{bin}} = f_s/N = 10$~Hz.
La resoluci\'on m\'inima en Hz se obtiene pasando $\Delta\omega$ (en rad/muestra)
a Hz:

\begin{center}
\begin{tabular}{lcc}
\toprule
Ventana & $\Delta\omega$ & $\Delta f_{\min}$ \\
\midrule
Rectangular     & $4\pi/N$   & $20$~Hz \\
Hann            & $8\pi/N$   & $40$~Hz \\
Hamming         & $8\pi/N$   & $40$~Hz \\
Blackman-Harris & $12\pi/N$  & $60$~Hz \\
Flattop         & $\approx 18\pi/N$ & $90$~Hz \\
\bottomrule
\end{tabular}
\end{center}

Para separar dos tonos a $50$~Hz:
sirven \textbf{Rectangular, Hann y Hamming} (las tres tienen $\Delta f_{\min} < 50$~Hz).
No sirven Blackman-Harris ni Flattop con este $N$.

\textbf{b) ¿Qu\'e ventanas no enmascaran?}
La regla simple: si el primer l\'obulo lateral est\'a por debajo de los $-40$~dB,
el pico d\'ebil ``sobresale'' por encima del piso de l\'obulos laterales.

\begin{itemize}
  \item Rectangular ($-13$~dB): \textbf{enmascara}.
  \item Hann ($-32$~dB): \textbf{enmascara} (apenas, pero s\'i).
  \item Hamming ($-43$~dB): \textbf{no enmascara}.
  \item Blackman-Harris ($-92$~dB): \textbf{no enmascara}.
  \item Flattop ($-93$~dB): \textbf{no enmascara}.
\end{itemize}

\textbf{c) Por qu\'e Flattop en metrolog\'ia.}
Lo especial de Flattop es que su l\'obulo principal es \emph{aproximadamente plano} en la cima
(de ah\'i el nombre).
Eso significa que si la frecuencia de la se\~nal no cae exactamente en un bin,
la altura del pico observado en el espectro es pr\'acticamente la misma que la
de la se\~nal real.
En cualquier otra ventana, si la frecuencia cae entre dos bins, la altura
medida es menor que la real (hasta $-3{,}9$~dB para la rectangular).

\textbf{Resumen:} Flattop sacrifica resoluci\'on para que la
\emph{amplitud medida} sea correcta sin importar d\'onde caiga la frecuencia.
Se usa en calibraci\'on de instrumentos y medici\'on de vibraciones.

% =========================================================
\section*{Ejercicio 3: Sesgo y varianza}

\begin{enun}
\textbf{a)} Defina con sus propias palabras:
\begin{itemize}[noitemsep]
  \item ?`Qu\'e es el \emph{sesgo} de un estimador?
  \item ?`Qu\'e es la \emph{varianza} de un estimador?
  \item ?`Qu\'e quiere decir que un estimador sea \emph{consistente}?
\end{itemize}

\textbf{b)} En la TS4 se estudia el estimador de amplitud
$\hat{a}_1 = |X_w^i(\Omega_0)|$ donde la frecuencia real $\Omega_1$ tiene
un peque\~no desv\'io aleatorio respecto a $\Omega_0$.

Para cada propiedad pedida, indique qu\'e ventana convendr\'ia usar
y por qu\'e:

\begin{center}
\begin{tabular}{lc}
\toprule
Quiero estimar... & Conviene la ventana... \\
\midrule
Amplitud con poco sesgo            & \rule{4cm}{0.3pt} \\
Frecuencia con poca varianza       & \rule{4cm}{0.3pt} \\
\bottomrule
\end{tabular}
\end{center}

\textbf{c)} ?`Existe una ventana que sea \'optima para ambas cosas?
?`Por qu\'e s\'i o por qu\'e no?
\end{enun}

\res

\textbf{a) Definiciones.}

\textbf{Sesgo:} es la diferencia entre lo que mi estimador devuelve
\emph{en promedio} y el valor verdadero del par\'ametro.
\[
  s = E\{\hat{\theta}\} - \theta
\]
Si repito el experimento muchas veces y promedio los resultados,
?`me acerco al valor real?
Si s\'i $\Rightarrow$ sesgo bajo.
Si me alejo siempre en el mismo sentido $\Rightarrow$ sesgo alto.
Es un error sistem\'atico.

\textbf{Varianza:} es cu\'anto var\'ian los resultados del estimador entre
realizaciones distintas.
\[
  v = E\{(\hat{\theta} - E\{\hat{\theta}\})^2\}
\]
Si repito el experimento muchas veces y los resultados son muy parecidos
entre s\'i $\Rightarrow$ varianza baja.
Si los resultados ``saltan'' mucho de una realizaci\'on a otra $\Rightarrow$
varianza alta. Es el error aleatorio.

\textbf{Consistente:} es un estimador que cumple las dos cosas a medida que
$N \to \infty$: el sesgo tiende a cero \emph{y} la varianza tiende a cero.
Es decir, con suficientes muestras, el estimador converge al valor real.

\textbf{Analog\'ia del tirador al blanco:} sesgo es si todos los tiros le pegan
arriba del blanco; varianza es si los tiros est\'an muy dispersos entre s\'i.
Un buen estimador tiene los tiros agrupados \emph{y} en el centro.

\textbf{b) Qu\'e ventana usar para cada caso.}

\begin{center}
\begin{tabular}{lll}
\toprule
Quiero estimar... & Ventana & Por qu\'e \\
\midrule
Amplitud con poco sesgo   & \textbf{Flattop}     & L\'obulo principal plano \\
Frecuencia con poca varianza & \textbf{Rectangular} & L\'obulo principal m\'as estrecho \\
\bottomrule
\end{tabular}
\end{center}

\emph{Amplitud:} Flattop tiene su l\'obulo principal aplanado, as\'i que el pico
medido se acerca al pico real sin importar d\'onde caiga la frecuencia.
Resultado: sesgo bajo en amplitud.

\emph{Frecuencia:} Rectangular tiene el l\'obulo principal m\'as estrecho,
o sea el pico m\'as ``puntiagudo''.
Localizar el m\'aximo es m\'as preciso, y entre realizaciones distintas
da resultados m\'as parecidos $\Rightarrow$ varianza baja en la
estimaci\'on de frecuencia.

\textbf{c) ?`Hay una ventana \'optima para ambas?}

No. Es un trade-off fundamental:
\begin{itemize}
  \item Para medir bien la \emph{amplitud} quiero un l\'obulo principal
    ancho y plano $\Rightarrow$ peor resoluci\'on en frecuencia.
  \item Para medir bien la \emph{frecuencia} quiero un l\'obulo principal
    estrecho y puntiagudo $\Rightarrow$ peor estimaci\'on de amplitud
    cuando la frecuencia no cae en un bin.
\end{itemize}

No existe una ventana que tenga un l\'obulo principal estrecho y al mismo
tiempo plano. Es un compromiso: hay que decidir qu\'e me importa m\'as
seg\'un la aplicaci\'on.

% =========================================================
\section*{Ejercicio 4: Tipo parcial --- Elecci\'on de ventana, $N$ y $f_s$}

\begin{enun}
Se tiene una se\~nal compuesta por:
\begin{itemize}[noitemsep]
  \item Una senoidal de \textbf{$0$~dB} en $f_1 = 1000$~Hz.
  \item Una cosenoidal de \textbf{$-40$~dB} en $f_2 = 1050$~Hz.
  \item Ruido gaussiano de fondo de $-60$~dB.
\end{itemize}
Quiero ver ambas se\~nales separadas en el espectro.

Complete la siguiente tabla decidiendo, para cada ventana, si sirve o no
para resolver el problema. Justifique brevemente.

\medskip
\begin{center}
\small
\begin{tabular}{lcccc}
\toprule
Ventana & $N_{\min}$ p/ resolver & ?`Detecta $-40$~dB? & ?`Sirve? \\
\midrule
Rectangular     & \rule{2cm}{0.3pt} & \rule{1cm}{0.3pt} & \rule{1.5cm}{0.3pt} \\
Hann            & \rule{2cm}{0.3pt} & \rule{1cm}{0.3pt} & \rule{1.5cm}{0.3pt} \\
Hamming         & \rule{2cm}{0.3pt} & \rule{1cm}{0.3pt} & \rule{1.5cm}{0.3pt} \\
Blackman-Harris & \rule{2cm}{0.3pt} & \rule{1cm}{0.3pt} & \rule{1.5cm}{0.3pt} \\
Flattop         & \rule{2cm}{0.3pt} & \rule{1cm}{0.3pt} & \rule{1.5cm}{0.3pt} \\
\bottomrule
\end{tabular}
\end{center}

Use $f_s = 8000$~Hz. La condici\'on de resoluci\'on es
$\Delta f_{\min} < 50$~Hz (separaci\'on entre las dos se\~nales).

\textbf{Tip:} para calcular $N_{\min}$, ponga
$\Delta\omega / (2\pi) \cdot f_s < 50$ y despeje $N$.
\end{enun}

\res

Para cada ventana, $N_{\min}$ sale de la condici\'on
$\Delta f_{\min} < 50$~Hz.
Si el primer l\'obulo lateral est\'a por debajo de $-40$~dB,
la ventana detecta la se\~nal d\'ebil.

\begin{center}
\begin{tabular}{lccc}
\toprule
Ventana & $N_{\min}$ p/ resolver & ?`Detecta $-40$~dB? & ?`Sirve? \\
\midrule
Rectangular     & 321  & No  & \textbf{No} (no detecta) \\
Hann            & 641  & No  & \textbf{No} (apenas) \\
Hamming         & 641  & S\'i & \textbf{S\'i} \\
Blackman-Harris & 961  & S\'i & S\'i (pero pide m\'as $N$) \\
Flattop         & 1441 & S\'i & S\'i (pide mucho $N$) \\
\bottomrule
\end{tabular}
\end{center}

\textbf{Ejemplo de c\'alculo para Hamming:}
\[
  \frac{8\pi/N}{2\pi}\cdot 8000 < 50
  \;\Rightarrow\; \frac{4 \cdot 8000}{N} < 50
  \;\Rightarrow\; N > 640
  \;\Rightarrow\; N_{\min} = 641
\]

\textbf{Ventana \'optima:} \textbf{Hamming}.
Es la que con menor $N$ cumple ambas condiciones.
Si se quiere usar potencia de 2 por la FFT, conviene $N = 1024$.
La resoluci\'on efectiva queda en $4 \cdot f_s/N = 4 \cdot 7{,}8 \approx 31$~Hz,
con margen sobre los $50$~Hz que pide la separaci\'on.

% =========================================================
\section*{Ejercicio 5: Periodograma, Bartlett y Welch}

\begin{enun}
En las clases vimos tres m\'etodos no param\'etricos para estimar la
densidad espectral de potencia (PSD) de una se\~nal:

\smallskip
\textbf{Periodograma:} $\hat{P}(e^{j\omega}) = \frac{1}{N}|X_N(e^{j\omega})|^2$.
Es la DFT al cuadrado, normalizada.

\textbf{Bartlett:} parte la se\~nal en $K$ bloques de longitud $L$
\emph{sin solapar}, calcula el periodograma de cada uno, y los promedia.

\textbf{Welch:} igual que Bartlett pero los bloques s\'i se solapan
(t\'ipicamente 50\%) y cada bloque se enventana antes de hacer la FFT.

\smallskip
\textbf{a)} ?`Por qu\'e el periodograma ``no es un buen estimador''?
Pista: pensar en sesgo y varianza, ?`mejoran al aumentar $N$?

\textbf{b)} ?`Por qu\'e Bartlett reduce la varianza? ?`Qu\'e sacrifica para
lograrlo?

\textbf{c)} ?`Por qu\'e Welch es mejor que Bartlett?
Mencione las dos mejoras que introduce.

\textbf{d)} (Tipo parcial)
Tengo $N = 1000$ muestras de una se\~nal aleatoria.
Quiero estimar la PSD con baja varianza, pero tambi\'en quiero resoluci\'on
espectral aceptable.
Decida una estrategia razonable: ?`cu\'antos bloques uso? ?`de qu\'e tama\~no?
?`qu\'e ventana? Justifique.
\end{enun}

\res

\textbf{a) Por qu\'e el periodograma no es un buen estimador.}
Tiene dos problemas:
\begin{itemize}
  \item Es \textbf{sesgado}. El valor esperado del periodograma es la PSD real
    convolucionada con una ventana de Bartlett (tri\'angulo).
    El sesgo desaparece cuando $N \to \infty$ (es asint\'oticamente insesgado),
    pero para $N$ finito siempre hay error sistem\'atico.
  \item Su \textbf{varianza no decrece con $N$}. Para ruido gaussiano,
    $\text{Var}\{\hat{P}_{\text{per}}\} \approx P_x^2(e^{j\omega})$, sin importar
    cu\'antas muestras tome.
\end{itemize}
Como la varianza no tiende a cero, el periodograma \textbf{no es un estimador
consistente}. Por eso aparecen los m\'etodos de Bartlett y Welch.

\textbf{b) Bartlett: c\'omo reduce varianza y qu\'e sacrifica.}

Al promediar $K$ periodogramas independientes, la varianza del promedio se
reduce por un factor $K$:
\[
  \text{Var}\{\hat{P}_B\} \approx \frac{1}{K} P_x^2(e^{j\omega})
\]

\textbf{Lo que sacrifica:}
si la se\~nal total tiene $N$ muestras y la parto en $K$ bloques, cada bloque
tiene solo $L = N/K$ muestras.
Bloques m\'as cortos $\Rightarrow$ peor resoluci\'on espectral
($\Delta\omega \propto 1/L$).

\textbf{Trade-off:} m\'as bloques $\Rightarrow$ menos varianza pero peor resoluci\'on.
Menos bloques $\Rightarrow$ mejor resoluci\'on pero m\'as varianza.

\textbf{c) Por qu\'e Welch es mejor que Bartlett.}

Welch agrega dos mejoras:

\emph{1.~Enventanado de cada bloque.}
Cada bloque se multiplica por una ventana antes de la FFT.
Esto reduce el sesgo respecto a Bartlett (que usa bloques rectangulares
con sus l\'obulos laterales feos).

\emph{2.~Solapamiento entre bloques (t\'ipicamente 50\%).}
Al solapar, entran m\'as bloques en el mismo $N$ total, lo que permite
promediar m\'as veces sin acortar tanto cada bloque.
Por eso Welch reduce la varianza casi a la mitad de Bartlett
para el mismo $N$ y $L$.

\textbf{Resultado:} Welch logra mejor balance entre sesgo, varianza y
resoluci\'on que Bartlett. Por eso es el m\'etodo est\'andar para
estimar PSD de procesos estoc\'asticos.

\textbf{d) Estrategia con $N = 1000$.}

Una elecci\'on razonable:
\begin{itemize}[noitemsep]
  \item Bloques de \textbf{$L = 256$} muestras (potencia de 2).
  \item Solapamiento del 50\% (cada bloque comparte 128 muestras con el siguiente).
  \item Cantidad de bloques: aproximadamente $K \approx 7$ (se calcula como
    $K = \lfloor (N - L)/(L/2) \rfloor + 1$).
  \item Ventana \textbf{Hann} (uso general, buen balance).
\end{itemize}

\textbf{Justificaci\'on:} bloques de 256 dan resoluci\'on
$\Delta f \approx f_s / 256$, suficiente para la mayor\'ia de las
se\~nales. Al promediar 7 bloques la varianza se reduce
aproximadamente $7\times$ respecto al periodograma simple.
Si la se\~nal es muy ruidosa y prefiero menos varianza, puedo usar $L = 128$
y promediar m\'as bloques (sacrificando resoluci\'on).

% =========================================================
\section*{Ejercicio 6: A jugar con La440.wav}

\begin{enun}
Se proveen los archivos \texttt{La440.wav} y \texttt{La440\_delay.wav}.
La idea es que pruebe distintas ventanas con se\~nales reales y vea c\'omo
cambia el espectro.

\textbf{a)} Cargue \texttt{La440.wav} y graf\'iquelo en el dominio del tiempo
y en frecuencia (con ventana rectangular).
?`Aparece el pico esperado en 440~Hz? ?`Aparecen arm\'onicos?

\textbf{b)} Pruebe el mismo an\'alisis con Hann y con Flattop.
?`Cu\'al de las tres ventanas le da el pico m\'as estrecho?
?`Cu\'al le da la amplitud m\'as estable (o sea, m\'as cerca de un valor
``redondo'')?

\textbf{c)} Simule una segunda se\~nal cercana sumando al audio una
senoidal de la misma amplitud pero a 450~Hz.
?`Se pueden distinguir las dos componentes con cada ventana?
?`Y si pone la segunda a 600~Hz?

\textbf{d)} Repita el an\'alisis con \texttt{La440\_delay.wav}.
?`Se nota la diferencia mirando el m\'odulo del espectro?
?`Por qu\'e?
\end{enun}

\res

\textbf{a) Espectro b\'asico.}
La nota La debe mostrar un pico en 440~Hz.
Si el audio proviene de un instrumento real (no un seno puro), tambi\'en se
ver\'an arm\'onicos en 880~Hz (segundo arm\'onico), 1320~Hz (tercero), etc.
La altura relativa de los arm\'onicos depende del timbre del instrumento.

\textbf{b) Comparando ventanas.}
\begin{itemize}
  \item \textbf{Rectangular}: pico m\'as estrecho, pero con ``patas'' de leakage si 440~Hz no cae en bin exacto.
  \item \textbf{Hann}: pico m\'as ancho, patas casi desaparecidas.
  \item \textbf{Flattop}: pico claramente m\'as ancho, pero la altura del pico
    es la m\'as confiable (no depende de si 440 cae en bin).
\end{itemize}

\textbf{Pico m\'as estrecho:} rectangular.
\textbf{Amplitud m\'as estable:} Flattop.

\textbf{c) Segunda senoidal cercana.}

Para hacer la prueba, basta con sumar al \texttt{x} cargado:
\begin{lstlisting}[style=pystyle]
n  = np.arange(len(x))
x2 = x + np.cos(2*np.pi * 450 * n / fs)  # cambiar 450 por 600 para la 2da prueba
\end{lstlisting}

\emph{Con 450~Hz} (separaci\'on $\Delta f = 10$~Hz):
\begin{itemize}[noitemsep]
  \item Si $N$ no es muy grande, todas las ventanas mostrar\'an un solo pico
    ensanchado, no dos.
  \item Solo con $N$ grande \emph{y} ventana rectangular se logra separar
    visualmente los dos picos.
\end{itemize}

\emph{Con 600~Hz} (separaci\'on $\Delta f = 160$~Hz, muy grande):
\begin{itemize}[noitemsep]
  \item Las dos componentes se separan f\'acilmente con cualquier ventana.
  \item Es buen momento para probar Flattop: como los picos est\'an lejos,
    el ancho extra de su l\'obulo principal no es un problema, y la
    estimaci\'on de amplitud queda muy estable.
\end{itemize}

\textbf{d) Por qu\'e el retardo no se ve en el m\'odulo.}

Un retardo temporal de $\tau$ muestras es, en frecuencia, un cambio de fase:
\[
  X_{\text{delay}}[k] = X[k] \cdot e^{-j 2\pi k \tau / N}
\]
El factor $e^{-j2\pi k\tau/N}$ tiene m\'odulo 1.
Por lo tanto $|X_{\text{delay}}[k]| = |X[k]|$: el espectro de m\'odulo es
\emph{id\'entico} en ambos archivos.

Para ver el retardo hay que mirar la \textbf{fase} del espectro
($\angle X[k]$), que cambia linealmente con $k$.
Otra opci\'on es la \textbf{correlaci\'on cruzada} entre ambas se\~nales,
que muestra un pico en el valor del retardo.

\end{document}
